\begin{tikzpicture}[
  font=\scriptsize,
  kotak/.style={draw, rounded corners, minimum height=0.95cm, minimum width=2.5cm,
                align=center, font=\scriptsize\bfseries},
  panah/.style={-{Stealth[length=2mm]}, thick},
  duaarah/.style={{Stealth[length=2mm]}-{Stealth[length=2mm]}, thick},
  ket/.style={font=\tiny, align=center, fill=white, inner sep=1.5pt}
]
  \node[kotak, fill=praditagreen!15, minimum width=2.7cm, minimum height=1.1cm] (inti)
    at (0,0) {Medikara+\\{\tiny\mdseries platform}};

  \node[kotak, fill=orange!12] (pasien) at (-4.6,0)
    {Pasien\\{\tiny\mdseries mencari dan memesan}};
  \node[kotak, fill=orange!12] (dokter) at (4.6,0)
    {Dokter klinik mitra\\{\tiny\mdseries membuka jadwal}};
  \node[kotak, fill=blue!8] (perusahaan) at (-2.4,2.0)
    {Perusahaan klien\\{\tiny\mdseries layanan karyawan}};
  \node[kotak, fill=blue!8] (asuransi) at (2.4,2.0)
    {Asuransi swasta\\{\tiny\mdseries penjaminan biaya}};
  \node[kotak, fill=blue!8, minimum width=3.4cm] (apotek) at (0,-2.75)
    {Apotek dan laboratorium mitra\\{\tiny\mdseries resep dan hasil pemeriksaan}};

  \draw[duaarah] (pasien) -- node[ket, above] {cari, pesan, bayar} (inti);
  \draw[duaarah] (dokter) -- node[ket, above] {jadwal, layanan} (inti);
  \draw[duaarah] (perusahaan) -- (inti);
  \draw[duaarah] (asuransi) -- (inti);
  \draw[duaarah] (apotek) -- (inti);

  % Busur dilandaikan agar tidak menyentuh kotak apotek, dan labelnya di luar panah tengah.
  \draw[panah, dashed, gray!70!black] (pasien.south) to[out=-40, in=-140, looseness=0.6]
    node[ket, below, pos=0.2, xshift=-4pt] {interaksi langsung:\\janji temu dan konsultasi}
    (dokter.south);

  \node[font=\tiny, anchor=north, align=center, text=gray!60!black] at (0,-3.4)
    {oranye: sisi interaksi inti \quad biru: calon sisi tambahan};
\end{tikzpicture}
